\documentclass[11pt]{article}
\usepackage[margin=1in]{geometry}
\usepackage{amsmath,amssymb,amsthm}
\usepackage[T1]{fontenc}
\usepackage{lmodern}
\usepackage[hidelinks]{hyperref}
\newtheorem{theorem}{Theorem}
\newtheorem{lemma}{Lemma}
\title{The normalization obstruction in prime-time averages\\Conjecture 00000001695}
\author{gaochengzhecpu}
\date{October 9, 2026}
\begin{document}
\maketitle
\begin{abstract}
Use the constant observable on the unit circle rotated by $\sqrt2$.
This is a measure-preserving linear polynomial phase with a badly
approximable coefficient. With the source's normalization by $N$, its
prime-time average stays at least $1/4$ away from its integral for every
$N\geq16$. The asserted rate fails for every nonnegative constant at
arbitrarily large $N$.
\end{abstract}

\section{An admissible rotation}
Let $X=\mathbb R/\mathbb Z$ with normalized Haar measure $\mu$, let
$T(x)=x+\sqrt2\pmod1$, and let $f(x)=1$. Then $T$ preserves $\mu$ and
$T^p(x)=x+p\sqrt2\pmod1$, a linear polynomial phase. In particular
$\int_X f\,d\mu=1$. The coefficient satisfies the following explicit
bad-approximation condition.

\begin{lemma}
For every $p\in\mathbb Z$ and integer $q\geq1$,
\[
 \left|\sqrt2-\frac pq\right|\geq\frac{1}{5q^2}.
\]
\end{lemma}
\begin{proof}
If the left side is at least one, the assertion follows from $q\geq1$.
Otherwise $|p/q|<\sqrt2+1<3$, so $|\sqrt2+p/q|<5$.
The integer $2q^2-p^2$ is nonzero: equality would make $\sqrt2$ rational.
Consequently
\[
 1\leq|2q^2-p^2|
 =q^2\left|\sqrt2-\frac pq\right|
       \left|\sqrt2+\frac pq\right|
 <5q^2\left|\sqrt2-\frac pq\right|.
\]
This gives the required inequality.
\end{proof}

\section{The prime-counting obstruction}
Write $\pi(N)=\#\{p\leq N:p\text{ is prime}\}$. The average in the source is
\[
 A_N(x)=\frac1N\sum_{p\leq N}f(T^p x)=\frac{\pi(N)}{N}.
\]
An elementary sieve using only the parity of primes gives the convenient
bound $2\pi(N)\leq N+4$ for $N\geq2$. Thus for $N\geq16$,
\[
 \frac{\pi(N)}N\leq\frac12+\frac2N\leq\frac34,
 \qquad
 \left|A_N(x)-\int_X f\,d\mu\right|\geq\frac14.
\]
No prime number theorem or quantitative distribution theorem is needed.

\newpage
\section{Every proposed constant fails}
\begin{theorem}
For every $C\geq0$, every natural threshold $N_0$, and every $x\in X$,
there is an integer $N\geq\max(N_0,16)$ such that
\[
 C\frac{\log\log N}{\sqrt N\log N}
 <\left|A_N(x)-\int_X f\,d\mu\right|.
\]
\end{theorem}
\begin{proof}
For $N\geq2$, $\log N>0$ and $\log(\log N)\leq\log N$. Therefore
\[
 C\frac{\log\log N}{\sqrt N\log N}\leq\frac C{\sqrt N}.
\]
Choose an integer $N>(4C+1)^2+N_0+16$. Then $N\geq N_0$, $N\geq16$,
and $\sqrt N>4C$, so $C/\sqrt N<1/4$. Apply the preceding uniform
lower bound on the error.
\end{proof}
Thus no nonnegative constant works even eventually. The observable is
smooth and bounded, so regularity cannot repair this estimate.

\paragraph{Scope.}
The counterexample uses the stated normalization $1/N$ and integral.
There is no zero-mean hypothesis. A $\pi(N)$ normalization or logarithmic
prime weights give different expressions, outside this disproof's scope.

\section{Formal verification}
Lean uses the actual \texttt{UnitAddCircle}, Haar volume, Bochner integral,
function iteration, prime finsets, square root, and logarithm.
\begingroup\raggedright
\begin{itemize}
\setlength{\itemsep}{3pt}
\item \texttt{sqrt\_\allowbreak two\_\allowbreak badly\_\allowbreak approximable}
proves the explicit $1/(5q^2)$ bound. The rotation preserves measure and its
iterates are the displayed linear phase.
\item \texttt{actual\_\allowbreak average} and \texttt{actual\_\allowbreak integral}
compute the prime sum and integral. \texttt{prime\_\allowbreak count\_\allowbreak bound}
uses the sieve with modulus two; \texttt{error\_\allowbreak lower\_\allowbreak bound}
proves the resulting uniform error bound.
\item \texttt{arbitrarily\_\allowbreak large\_\allowbreak failure} proves the theorem for
every $C\geq0$, threshold, and point. The final \texttt{counterexample}
combines it with the actual rotation and bad-approximation property.
\end{itemize}
\endgroup
Only standard logical axioms are used, with no unproved placeholders,
custom axioms, native evaluation, or numerical experiments.

\paragraph{Reproduction.}
In \texttt{lean/}, run:
\begin{verbatim}
lake build
lake env lean -DwarningAsError=true Main.lean
\end{verbatim}
Lean 4.19.0 and an exact Mathlib revision are pinned. Run
\texttt{tectonic main.tex} to rebuild this PDF.

\paragraph{Sources.}
\href{https://github.com/The-Last-Math-Competition/The-Last-Math-Competition/blob/9b795e7a94a6076e49a65a6489e1caf9153abd23/conjectures/00000001695.md}{Conjecture 00000001695};
\href{https://github.com/leanprover-community/mathlib4/commit/c44e0c8ee63ca166450922a373c7409c5d26b00b}{Mathlib c44e0c8ee63c}.
\end{document}
